\section{Paso 9: análisis de los datos de salida}
\label{sec:paso9-analisis}
% Origen: analisis_corridas.py y resultados_informe.py sobre corridas_peatonales.csv (tablas en ../generado/).
% Variación entre días, secundarias, efecto de los NAC y n*: anexo G.

\textbf{Método.} Se aplica lo fijado antes de correr (sección~\ref{sec:paso7-diseno-experimentos}): para cada
alternativa y medida se forma la diferencia apareada $Z_j$ y se construye el intervalo $t$ sobre las $n$
diferencias \parencite{law2015}, con nivel de Bonferroni $1 - 0{,}05/9$ para las primarias. Con $n = 6$ pares,
$t_{5;\,1-0{,}05/18} = 4{,}655$.

\subsection{Resultados por escenario}

La tabla~\ref{tab:resultados-escenarios} muestra la media de cada medida entre las réplicas.

\entradaGenerada{tabla-resultados}

\textbf{Niveles absolutos.} En la situación actual (\EO{}), el pasajero típico espera unos 8~s en la cola. El
P90 es de 19,8~s y el 5,4~\% espera más de 30~s. En el pico de 08:15 a 08:45 el P90 sube a 24,6~s. Son esperas
cortas en promedio, coherentes con un sistema subutilizado: la utilización media de los molinetes es del
24~\%. La cola existe, pero es transitoria: en promedio hay 15,5 personas en cola y el máximo del día llega
a 64. El P90 de \EO{} va de 14,7 a 29,7~s según el día, pero el orden de los escenarios se
repite los seis días: $\EIT{} < \EI{} \approx \ET{} < \EO{}$. Las secundarias bajan en el mismo orden; el
tránsito del Roca y el tiempo de drenaje no sirven para comparar escenarios (anexo~\ref{anx:variacion}).

\subsection{Comparación con la situación actual}

La tabla~\ref{tab:diferencias} y la figura~\ref{fig:forest} presentan las diferencias de cada alternativa
contra \EO{} en las tres medidas primarias, con su intervalo de Bonferroni.

\entradaGenerada{tabla-diferencias}

\begin{figure}[H]
\centering
\figuraSiExiste{tpi-subte-res-forest.png}{0.85\textwidth}
\caption{Diferencia de cada alternativa contra \EO{} en las medidas primarias, con su intervalo de confianza
simultáneo (Bonferroni, $k = 9$). Un intervalo entero a la izquierda del cero es una mejora con evidencia
estadística.}
\label{fig:forest}
\end{figure}

\textbf{Regla de decisión.} Los nueve intervalos quedan enteros por debajo de cero. Con 6 días y confianza
conjunta de al menos el 95~\%, cada alternativa mejora las tres medidas primarias frente a \EO{}. \EI{}
(22 molinetes) baja el P90 en 7,95~s ($-40$~\%), la proporción con espera mayor a 30~s en 4,0 puntos
($-74$~\%) y el P90 del pico en 11,25~s ($-46$~\%); \ET{} (validación rápida), en 7,35~s ($-37$~\%), 2,9 puntos
($-53$~\%) y 9,70~s ($-40$~\%), y \EIT{} (las dos), en 11,25~s ($-57$~\%), 5,0 puntos ($-92$~\%) y 15,20~s
($-62$~\%). Los intervalos son anchos: el de \EI{} en el P90 va de $-15{,}8$ a $-0{,}1$~s. Lo que está
establecido es el signo de la mejora; su tamaño, con 6 días, queda impreciso.

\textbf{Entre alternativas.} La regla pide recomendar la de menor P90 solo si el intervalo frente a las otras
también lo sostiene. Fuera de la familia principal, con el mismo nivel de 99,44~\%, \EIT{} tiene menor P90 que
\EI{} ($-3{,}3$~s, IC $[-4{,}6;\,-2{,}0]$) y que \ET{} ($-3{,}9$~s, IC $[-6{,}8;\,-1{,}0]$), y entre \EI{} y
\ET{} no hay diferencia demostrable ($-0{,}6$~s, IC $[-2{,}4;\,1{,}2]$). En la proporción mayor a 30~s y en el
P90 del pico, los intervalos de \EIT{} contra \ET{} contienen al cero: la ventaja de \EIT{} se sostiene en el
P90 de la franja, que es la medida principal.

\textbf{Números aleatorios comunes y réplicas.} La correlación del P90 entre \EO{} y cada alternativa es de
0,98, y la varianza de la diferencia queda entre un 43 y un 76~\% por debajo de la de corridas independientes:
sin números comunes, los intervalos habrían sido entre 1,3 y 2 veces más anchos (anexo~\ref{anx:nac}). Con los
semianchos fijados, $n^{*}$ va de 87 a 338 días para los percentiles y de 20 a 54 para la proporción
(anexo~\ref{anx:n-necesario}), cifras inalcanzables. Se decidió no agregar días: la pregunta es si cada
alternativa mejora, y eso lo responde el signo de los nueve intervalos, que ya excluyen al cero. El tamaño de la
mejora queda impreciso y se informa con su intervalo. Con $n = 6$, por debajo del $n_0 = 10$ planificado, el
resultado se apoya en un $t$ crítico más exigente (4,655 con 5 grados de libertad, contra 3,622 con 9).
